\documentclass[11pt,portuguese,a4paper]{article}
           
%\usepackage{amssymb}
%\usepackage{amsmath,amsfonts,amsthm,amstext}
\usepackage[brazil]{babel}
\usepackage[latin1,utf8]{inputenc}
\usepackage{fancyheadings}
%\usepackage[dvips]{graphicx}
\usepackage{graphicx}
%\usepackage{latexsym}
%\usepackage{psfrag}

\setlength{\paperheight}{297mm} %definição do tamanho do papel - A4
\setlength{\paperwidth}{210mm} %definição do tamanho do papel - A4
\setlength{\oddsidemargin}{0.0mm} %coloca margem a 30.0mm 
\setlength{\evensidemargin}{-5.4mm} %coloca margem a 20.0mm 
\setlength{\topmargin}{-5.4mm} %coloca margem a 25.0mm 
\setlength{\headheight}{7mm} %espaço reservado para cabeçalho nas páginas
\setlength{\headsep}{7mm} %espaço entre o cabeçalho e o corpo do texto
\setlength{\textwidth}{160mm} %largura do corpo do texto -> 25mm de margens
\setlength{\textheight}{235mm} %altura do corpo do texto  -> 25mm de margens
\setlength{\parindent}{10mm} %tamanho do parágrafo

\pagestyle{fancy}
 \lhead{\includegraphics[scale=0.975]{letterhead-new}}
 \chead{}
 \rhead{}

 \lfoot{}
 \cfoot{}
 \rfoot{}

\setlength{\headrulewidth}{0.pt}
\setlength{\footrulewidth}{0.pt}


\begin{document}


\begin{center}

$\qquad$

\vspace{12pt}

\textsc{{\LARGE {Sistema de Informática do MAT/UnB}}}

\vspace{12pt}

{\Large Abertura de nova conta em LATEX_DATAABERTURA para}

\vspace{12pt}

{\LARGE \it LATEX_NOME LATEX_SOBRENOME }

\end{center}

\vspace{12pt}

\noindent {\Large \underline{Dados da conta:}}\\

\noindent {\LARGE {\bf Login}: LATEX_LOGIN }\\

\noindent {\LARGE {\bf Senha}: LATEX_SENHA  } \\

\noindent {\LARGE {\bf Email}: LATEX_MAIL } \\

\noindent {\LARGE {\bf Data de expiração da conta}: LATEX_DATAENCERRAMENTO } \\

\noindent {\small {\it Atenção}: Seu login e sua senha são sensíveis a letras maiúsculas e minúsculas. Por outro lado, seu email não é sensível a letras maiúsculas e minúsculas. Caso seja necessário manter sua conta ativa por mais tempo, entre em contato com o administrador do sistema através do email {\it redeadm@mat.unb.br}.}

\vspace{12pt}

\noindent {\Large \underline{Dados cadastrais:}}\\

\noindent {\large {\bf Nome completo do usuário}: LATEX_NOME LATEX_SOBRENOME }\\
\noindent {\large {\bf Email pessoal para contato}: LATEX_OUTROSMAILS }\\
\noindent {\large {\bf Vínculo com o MAT/UnB}: LATEX_VINCULO }\\
\noindent {\large {\bf Matrícula}: LATEX_MATRICULA }\\

\noindent {\small {\it Atenção}: Lembre-se de informar qualquer alteração nestes dados à secretaria do MAT/UnB ou pelo email {\it redeadm@mat.unb.br} . É essencial que o MAT/UnB tenha seus dados cadastrais atualizados.}

\vspace{12pt}

\noindent {\Large \underline{Quotas da conta:}}\\

\noindent {\large {\bf Quota de impressão/cópias mensal}: LATEX_QUOTAIMPRESSAO }\\
\noindent {\large {\bf Quota de email}: LATEX_QUOTAMAIL }\\
\noindent {\large {\bf Quota de home}: LATEX_QUOTADISCO }

\vspace{12pt}

\noindent {\Large \underline{Sobre o sistema do MAT/UnB:}}\\

\small 

\vspace{-12pt}

\begin{itemize}

\item Esta conta lhe dá acesso ao sistema de informática do MAT/UnB. Você pode se conectar a qualquer um dos terminais de acesso localizados nas salas dos alunos de doutorado, nos laboratórios de informática e na sala de leitura. Para maiores detalhes sobre o sistema do MAT/UnB e suas funcionalidades (backup de dados, conexão de máquinas pessoais, servidores de cálculo, acesso remoto, etc), visite {\it http://www.mat.unb.br/informatica}. 

\item O sistema do MAT/UnB adota o padrão linux (Debian) e há uma variedade de programas instalados no sistema. Você pode apreder mais sobre o sistema no endereço {\it http://www.debian.org/doc/}. 

\end{itemize}

\begin{flushright}
{\it Seja bem-vindo e aproveite sua estadia no MAT/UnB!}
\end{flushright}

\end{document}

